\chapter{Discussion réflexive}
\label{ch:discussion}
% BUDGET : 4 pages.

\section{Ce que le modèle fait au territoire qu'il représente}
\arediger{la performativité du modèle}

\section{Porter le modèle de ses propres encadrants}
% La parite comme contrainte epistemique : fidelite a un artefact plutot qu'a une
% realite. Ce que cela protege, ce que cela empeche.
\arediger{la posture d'ingénieur-modélisateur}

\section{Positionnement : ce que je n'ai pas vu}
% Six mois, extériorite, aucune enquete de terrain, des donnees produites ailleurs et
% pour autre chose. Le cout mesurable de cette position : les 25 variantes Karusmart
% vides, le volet cantines differe. Ce n'est pas une clause de style.
\arediger{positionnement en Guadeloupe}

\section{L'honnêteté épistémique comme livrable}
% Pourquoi le catalogue des pieges est le vrai produit du stage, et pas un aveu.
\arediger{la vigilance documentée comme résultat}

\section{Limites et transférabilité}
\arediger{limites et portage à un autre territoire}
